% NOTE (migration 23 Nov 2026): Text moved unchanged from main.tex. Much of it was drafted with AI assistance; rewrite it in your own words before submission (the supervisor runs AI checks).
\chapter{Related Work}\label{literature-review}

\section{Purpose and review approach}
This chapter examines university timetabling from an educational and administrative perspective. The main questions are what problem each paper addresses, who benefits, what capabilities were implemented or studied, which educational requirements were considered, what the evaluation establishes, and what remains missing. Algorithms are mentioned only where necessary to understand a capability or a limitation.

The review is a focused reading of five retained timetabling papers and one additional paper on natural-language preference interpretation, not a systematic review of all timetabling research. Two concern examination timetabling, one concerns a course-scheduling tool, one analyses real-world course-timetabling benchmarks, and one surveys university timetabling. These are different kinds of evidence: a review chapter does not demonstrate a single deployed product, a benchmark study does not constitute a user study, and a prototype demonstration does not establish improved academic outcomes. The discussion therefore distinguishes implemented functionality, author-reported scheduling results, intended educational benefits, and future work. Page references below refer to the supplied PDF page order where journal pagination is unavailable. Incomplete publication metadata is retained as provisional rather than guessed.

\section{Course planning and access to required teaching}
\subsection{Sobh, Cousens and Patel: Course Scheduler}
\textbf{Educational problem and beneficiaries.} Sobh, Cousens and Patel~\cite{sobh_course_scheduler} address the difficulty of deciding which courses to offer and when to offer them so that students can meet degree requirements. Their motivating example is a senior student with only a small set of remaining courses: an unsuitable offering or a clash can obstruct completion of the degree. The primary operator is the department chairperson or administrator; instructors contribute availability and teaching preferences, while students benefit indirectly through access to suitable course combinations (supplied PDF, pp.~1--2).

\textbf{Implemented capabilities and educational requirements.} The described suite combines student course requirements with instructor teaching capabilities and load limits to identify possible course offerings. Instructors select teaching times and rank possible course sets. Administrators inspect schedule viability, view specific student conflicts, and adjust course times through the interfaces. The tool gives feedback on which student requirements a proposed schedule satisfies. Thus, its educational emphasis is course access and progression, rather than merely filling a weekly grid (pp.~12--17).

\textbf{Evidence and limits.} The paper illustrates the interfaces and shows example combinations with 100\% and 82\% viability, including a conflict explanation for an individual student. These are demonstration results, not a controlled evaluation of graduation rates, student satisfaction, or staff time savings. Helping students graduate sooner is a motivation, not a measured outcome. The authors explicitly state that the suite does not address classroom allocation or class-size issues. They also identify fixed time slots and dependence on instructors' chosen times as limitations; generating times from student needs is future work (pp.~16--17).

\textbf{Implications for this project.} A useful timetable must protect access to required teaching and explain conflicts in terms an administrator can act on. This supports conflict reporting, review, and safe manual adjustment in the proposed project. However, the proposed prototype starts with predefined course activities and student groups. It does not select semester course offerings, optimise degree completion, or provide individual enrolment-based schedules unless those extensions are approved.

\subsection{Müller, Rudová and Müllerová: real-world course timetabling}
\textbf{Educational problem and beneficiaries.} Müller et al.~\cite{muller2025itc} examine the diversity of real university course-timetabling needs through ITC~2019. The central educational issue is whether students can attend the required combination of lectures, seminars, and laboratories. Beneficiaries include students seeking compatible class combinations and administrators coordinating teaching across faculties, buildings, and institutional calendars.

\textbf{Capabilities represented in the study.} The paper describes thirty benchmark instances from ten universities using UniTime in practice. The competition model combines class time and room assignment with student sectioning: students must be assigned the appropriate class types and compatible lecture--seminar or lecture--laboratory combinations. It represents room capacities and unavailable periods, repeated meetings and week patterns, travel time between rooms, and requirements on gaps, breaks, teaching days, and daily load (journal pp.~248--250). A student conflict can arise from either overlapping classes or insufficient travel time. These are requirements represented by the benchmark and contributing institutions; the paper is not presenting a newly tested interface for all these functions.

\textbf{Educational significance.} The institution descriptions show why a single definition of timetable quality is inadequate. Examples include travel between sports facilities, joint programmes creating unusual course combinations, and concentrated attendance days for students combining study with work. Scheduling therefore supports participation in a programme under practical attendance conditions, while institutional priorities determine acceptable trade-offs.

\textbf{Evidence and limits.} The paper analyses instance characteristics and compares scheduling solutions. This provides substantial evidence that real timetabling requires more than avoiding simultaneous room use. Its evaluation concerns feasibility and scheduling penalties, including student conflicts; it does not establish improvements in learning, retention, or satisfaction. Competition results also cannot be transferred directly to a smaller prototype with a different scope.

\textbf{Implications for this project.} The proposed system should make its supported educational requirements explicit and assess them against the approved institutional policy. Initial group-based conflict checking is a deliberate simplification of individual sectioning. Travel-time checking, irregular week patterns, and student sectioning should not be claimed as implemented features of the prototype without corresponding requirements and tests.

\section[Examination access and workload]{Examination access, workload and institutional coordination}
\subsection{Abdennadher and Edward: examination timetabling at GUC}
\textbf{Educational problem and beneficiaries.} Abdennadher and Edward~\cite{abdennadher_guc_exams} describe examination scheduling at the German University in Cairo. At the time described, manual course scheduling reportedly took about three weeks and examination scheduling about one week. Shared courses and resources across faculties made separate faculty schedules unsuitable. Examination administrators are the direct users; students benefit from schedules that allow them to sit their examinations, including examinations retaken from earlier semesters (supplied PDF, pp.~1--3).

\textbf{Implemented capabilities and educational requirements.} The system assigns examinations to periods, distributes students among rooms and seats, and displays results through a graphical interface. Its requirements include avoiding simultaneous exams for the same group, preventing semester clashes within a major to accommodate retakes, allowing no more than one core exam per day, limiting daily examinations, separating difficult examinations, respecting available seats, and supporting fixed and unavailable periods. Difficulty ratings are described as supplied by lecturers and students. These features connect scheduling to assessment workload and access; they should not be interpreted as evidence that the software itself determines pedagogically appropriate difficulty ratings (pp.~3--7).

\textbf{Evidence and limits.} The authors report a test using Spring~2007 final-examination data: 140 courses, 5,500 students and 1,300 available seats, with a solution generated within a few seconds. This demonstrates scheduling functionality on the stated data. The historical manual preparation time is not a measured before-and-after comparison of the entire administrative workflow. The paper does not report effects on exam performance, stress, or satisfaction. Integration with the invigilation system and further timetable repair and optimisation are future work (pp.~7--8).

\textbf{Implications for this project.} The paper supports clear distinctions between student access rules, assessment-spacing rules, and resource restrictions. It also motivates collecting local institutional requirements rather than copying another university's rules. Examination spacing and retake accommodation are relevant if an examination extension is approved; they are not automatically requirements for the current weekly course-timetabling scope.

\subsection{Kaya: multi-campus examination decision support}
\textbf{Educational problem and beneficiaries.} Kaya~\cite{kaya2026drl} addresses large-scale examination scheduling at a multi-campus Turkish university. Examinations have different durations, and many cohorts require more seats than a single room provides. Coordinators need to organise examinations within campus and facility restrictions while maintaining compatible assessment schedules. Students and teaching staff are affected by the resulting placements, although coordinators are the direct decision-support users.

\textbf{Implemented capabilities and educational requirements.} The reported framework handles duration-based placements, allocation across multiple rooms, room types and capacities, campus restrictions, instructor availability, and cohort non-overlap. Its institutional decision-support workflow lets coordinators reopen selected examinations and revise placements while preserving locked assignments. The paper describes selecting courses, moving placements, accepting suggestions, and approving a schedule (Sections~4.7 and 5.11--5.12). These capabilities matter because institutional timetables are revised and reviewed rather than simply generated once.

\textbf{Evidence and limits.} The main comparison uses 1,127 representative examinations, 110 rooms and 520 instructors, derived from 1,332 raw events. The paper reports high assignment coverage and an independent audit finding no hard-rule violation among assigned examinations. Coverage and validity must be read together: unassigned exams still require resolution. Crucially, cohort conflicts were made active through synthetic conflict injection because naturally occurring cohort overlap was absent from the evaluated production period; an accessibility requirement was also inactive. The operational logs establish execution at institutional scale, but the forty recorded coordinator interventions were a controlled replay, not observed live coordinator behaviour. The study is limited to one institution and period, and it does not establish reduced student stress, improved marks, or live administrative time savings (Sections~5.12 and 6.3).

\textbf{Implications for this project.} This paper supports reviewed changes, preservation of approved decisions, explicit reporting of unscheduled activities, and independent checking after edits. It also demonstrates why an educational review must ask whether relevant requirements were exercised in realistic data. The proposed evaluation will distinguish operational records, simulated cases, and participant evidence rather than treating them as interchangeable.

\section{University timetabling as an institutional decision}
\subsection{Petrovic and Burke: University Timetabling}
\textbf{Problem and beneficiaries.} Petrovic and Burke~\cite{petrovic2004university} review university course and examination timetabling as a major administrative activity. Their chapter concerns the interests of timetable officers, students and teachers, particularly the difficulty of satisfying different institutional priorities within limited resources (supplied PDF, pp.~2--3).

\textbf{Capabilities discussed.} The chapter reviews approaches to constructing feasible schedules, balancing several quality criteria, reusing previous scheduling experience, and adapting methods across problems. Its educational contribution to this review is the distinction between mandatory attendance and resource rules and desirable preferences, alongside recognition that those rules vary across universities. It explains why course and examination scheduling need different assumptions about repeated meetings and room use. Because this is a review chapter, these capabilities should not be presented as features implemented together in a single system.

\textbf{Evidence and limits.} The reviewed studies provide scheduling results and examples of ways to represent competing objectives. The chapter does not supply one controlled evaluation of an integrated product or evidence of improved learning outcomes. Its range of examples supports context-sensitive requirements, not a universal checklist that every university should adopt.

\textbf{Implications for this project.} The proposed prototype should let administrators express supported local requirements, identify which are mandatory, and inspect which preferences remain unmet. The distinction is important educationally: avoiding a clash protects access, while satisfying a preferred teaching time improves convenience under a particular institutional policy. Neither should be conflated with a demonstrated change in educational achievement.

\section{Natural-language preferences and faithful interpretation}
\subsection{Sanguinetti et al.: From User Preferences to Optimization Constraints Using Large Language Models}
\textbf{Problem and beneficiaries.} Sanguinetti et al.~\cite{sanguinetti2025preferences} investigate how people can express household appliance needs in ordinary language so that those needs become constraints for an energy optimizer. The motivation is that a community's energy-saving objectives may conflict with individual routines: a technically efficient recommendation is of limited use if it ignores when a person needs hot water or air conditioning. The direct beneficiaries are household users in Italian renewable energy communities. Although this is not an educational study, the same concern motivates the proposed timetable interface: an administrator's intended attendance or availability requirement must survive its translation into a formal rule (supplied PDF, pp.~1--3).

\textbf{Capabilities studied.} The evaluated component converts Italian utterances into representations of appliance operating times or desired temperatures. It supports requests applying throughout the time horizon, within a specified interval, after a start time, or until an end time. The input includes both the utterance and manually identified preference spans marked with XML tags. The study therefore evaluates conversion of identified preferences, rather than the complete discovery of all requirements in unrestricted text. A connection to an external optimizer is described within the wider project; the paper does not establish a completed, evaluated workflow with user clarification, interpretation approval and downstream schedule validation (pp.~3--8).

\textbf{Evidence.} The authors manually annotate a pilot subset of 26 utterances drawn from a 157-utterance Italian corpus. Prompts provide no example, one example, or five examples. Evaluation separates character-level similarity (ChrF) from correctness of variables and time conditions. In the five-example setting, Table~1 reports variable accuracy of 0.8571 and condition accuracy of 0.4286 for LLaMAntino-3-ANITA; the corresponding ChatGPT values are 0.7619 and 0.3571. These are the paper's component-level measures, not the proportion of complete requests interpreted correctly. Conditions remain substantially harder than variable identification, and adding examples does not improve every component consistently. The authors conclude that the tested prompting approach is inadequate for this task and propose further data collection and supervised fine-tuning (pp.~5--10).

\textbf{Limits of the evidence.} The small, handcrafted sample, Italian appliance domain and supplied preference spans limit transfer to English university requirements containing multiple rules, ambiguous staff names or institutional terms. Character similarity cannot establish that a constraint preserves user intent. The supplied version also has a reporting inconsistency: its model-selection and results narrative describes four models and excludes two, while Table~1 additionally lists Cerbero, Maestrale and IT5 without matching experimental descriptions. Those additional rows are not used here to infer a reproducible ranking. The paper does not report a controlled comparison of manual and assisted entry, measured user effort, educational outcomes, or an evaluated clarification-and-approval process. Functional evaluation after connection to the optimizer is proposed as future work (pp.~7--10).

\textbf{Implications for this project.} This is a close conceptual precedent for converting everyday requirements into a restricted constraint vocabulary, but it does not establish that the proposed university application will work. The thesis will consequently assess complete semantic interpretation and individual fields separately, including time boundaries, entity references and mandatory versus preferred status. Requirements will be evaluated from the full statement, with the supplied institutional context; manually highlighted spans will not be assumed in the main evaluation. Review and targeted clarification are proposed safeguards whose effectiveness must be measured. The human comparison will assess final interpretation correctness alongside time and corrections, rather than treating a fluent response or a valid timetable as evidence that the administrator's intent was understood. These implications support the existing SRS requirements for entity resolution, named periods, clarification, review and evaluation (FR-016--FR-020, FR-049 and FR-050); they add no appliance-related functions to the project.

\section[Cross-paper synthesis]{Cross-paper synthesis and implications for the proposed system}
\textbf{Access and progression.} Sobh et al. emphasise course requirements and degree progression, Müller et al. emphasise attendance across compatible class combinations, and the GUC study explicitly considers retake students. Together, they show that avoiding clashes has educational meaning: students must be able to attend required teaching or assessment. However, offering selection, individual sectioning and group-based conflict prevention are different capabilities. The proposed prototype initially addresses the last of these.

\textbf{Workload and assessment organisation.} GUC's core-exam and difficulty-spacing requirements show how assessment policy affects scheduling. These rules are institution-specific; they do not establish an optimal universal spacing policy. Course gaps, late sessions and daily-load balance are related concerns for the proposed course prototype, but their thresholds and priority require local agreement.

\textbf{Staff coordination and usable decisions.} Course Scheduler's conflict feedback and Kaya's reviewed revisions show the importance of giving administrators actionable information and control. A useful system should expose incomplete schedules, violated preferences and the consequences of a change. Faster generation alone is insufficient if the administrator cannot understand or approve the result.

\textbf{Evidence quality.} Across the papers, most evaluation concerns scheduling feasibility, coverage, runtime or penalties. Those outcomes are relevant to institutional operations, but they do not directly measure learning, retention, graduation, fairness, stress or satisfaction. The thesis will therefore describe improved access and workload organisation as intended educational benefits, and claim improvements only for outcomes actually measured.

\textbf{Proposed functional response.} The project will support entry of institutional requirements in ordinary language, display the proposed interpretation, request clarification where needed, and require administrator approval. It will generate course schedules from approved rules; identify group, instructor and room conflicts; check capacity and facility suitability; support comparison and reviewed manual changes; and release validated timetables to authorised users. These are proposed capabilities linked to the SRS, not claims that a completed application already exists. Exam-specific and individual-enrolment features remain extensions.

\section{Research gap and limits of the review}
Within the five retained timetabling papers, the educational needs are represented as predefined rules, input data, preferences or choices made through interfaces. None evaluates the proposed combination of ordinary-language requirement entry, explicit interpretation review, useful clarification, and separate measurement of interpretation accuracy for university course scheduling. This bounded observation motivates studying whether administrators can record their scheduling needs with less effort while preserving their intended meaning. It is not a claim that no related system exists anywhere. Within the sources examined, Arabic and Egyptian Arabic administrative constraint extraction is not evaluated; the Cairo examination context does not itself establish natural-language input support. This motivates language-stratified evaluation without asserting broader novelty.

The additional study by Sanguinetti et al.~\cite{sanguinetti2025preferences} directly evaluates preference-to-constraint conversion in another domain, so the translation task itself is not claimed as new. It does not evaluate the proposed university review workflow. Other sources cited in the introduction provide further context rather than evidence of educational effectiveness in this project. A documented search of further educational timetabling and human-centred scheduling literature is needed before making a broader novelty claim. The search should record sources, dates, inclusion criteria and screening decisions, and include terms concerning student access, administrative workload, participatory scheduling and requirement entry alongside natural-language scheduling.

\section{Summary}
The reviewed literature identifies three priorities for the proposed study: protecting access to required teaching, expressing institution-specific scheduling needs, and keeping administrators informed and in control. Existing evidence mainly establishes scheduling capabilities and operational results. The methodology therefore evaluates both the correctness of recorded requirements and the effort and understanding involved in the administrative workflow, without claiming unmeasured improvements in student learning.
